\documentclass[12pt]{article}
\usepackage[T1]{fontenc}
\usepackage[utf8]{inputenc}
\usepackage{lmodern}
\usepackage{textcomp}
\usepackage{enumitem}
\usepackage{geometry}
\geometry{margin=1in}
\begin{document}
\begin{center}
Answer Key - History - Undergraduate Level Assessment 15 \
\small (Answers are ordered exactly as in the question paper.)
\end{center}

\section*{Multiple Choice}
\begin{enumerate}[label=\arabic*.]
\item \textbf{Answer:} B
\\ \textit{Options:} A) the emergence of complex social hierarchies and political systems; B) a growth in cultural diversity; C) the formation of a single global civilization; D) the onset of a new ice age; E) the introduction of farming from neighboring continents
\\ \textit{Note:} During the Holocene epoch, Australia, like many other regions worldwide, experienced significant social and cultural transformations, leading to increased cultural diversity as populations adapted to changing environments and developed distinct practices and identities.
\item \textbf{Answer:} A
\\ \textit{Options:} A) Mogollon; B) Navajo; C) Clovis; D) Paleo-Indians; E) Hopewell
\\ \textit{Note:} The Mogollon culture is recognized for its sedentary lifestyle in the Southwest due to their reliance on agriculture, particularly the cultivation of crops like maize, which necessitated permanent settlements and the development of complex societies.
\item \textbf{Answer:} A
\\ \textit{Options:} A) gas chromatography-mass spectrometry and atomic force microscopy; B) infrared spectroscopy and ultraviolet-visible spectroscopy; C) Raman spectroscopy and Fourier-transform infrared spectroscopy; D) electron microscopy and nuclear magnetic resonance; E) gravitational wave detection and atomic absorption spectroscopy
\\ \textit{Note:} Gas chromatography-mass spectrometry and atomic force microscopy are correct because both techniques are advanced analytical methods used for the precise detection and characterization of trace elements in various samples.
\item \textbf{Answer:} C
\\ \textit{Options:} A) construction of permanent dwellings; B) invention of the wheel; C) parietal and mobiliary art; D) creation of writing systems; E) development of pottery
\\ \textit{Note:} The presence of parietal and mobiliary art during the Upper Paleolithic strongly indicates advances in human intellect as it reflects complex cognitive abilities, symbolic thinking, and the capacity for abstract representation, which are hallmarks of advanced cultural development.
\item \textbf{Answer:} A
\\ \textit{Options:} A) China.; B) Africa.; C) Australia.; D) North America.; E) South America.
\\ \textit{Note:} The earliest-known use of bronze is attributed to China due to archaeological evidence from the Shang Dynasty, which dates back to around 1600-1046 BCE, demonstrating advanced metallurgy and the production of bronze artifacts.
\end{enumerate}

\section*{True / False}
\begin{enumerate}[label=\arabic*.]
\setcounter{enumi}{5}
\item \textbf{Answer:} False
\\ \textit{Options:} A) True; B) False
\\ \textit{Note:} The correct answer is: 1942
\item \textbf{Answer:} True
\\ \textit{Options:} A) True; B) False
\\ \textit{Note:} According to the passage: coal tar
\item \textbf{Answer:} False
\\ \textit{Options:} A) True; B) False
\\ \textit{Note:} The correct answer is: Scarlet
\item \textbf{Answer:} False
\\ \textit{Options:} A) True; B) False
\\ \textit{Note:} The correct answer is: endoheretics
\item \textbf{Answer:} False
\\ \textit{Options:} A) True; B) False
\\ \textit{Note:} The correct answer is: comics
\end{enumerate}

\section*{Long Form Response}
\begin{enumerate}[label=\arabic*.]
\setcounter{enumi}{10}
\item \textbf{Answer:} Ángel Trías
\\ \textit{Note:} Expected answer: Ángel Trías
\item \textbf{Answer:} 81 BC
\\ \textit{Note:} Expected answer: 81 BC
\end{enumerate}

\vfill
\noindent\textit{End of Answer Key}
\end{document}